\documentclass[11pt]{article}

\usepackage[margin=1in]{geometry}
\usepackage{amsmath,amssymb}
\usepackage{booktabs}
\usepackage{graphicx}
\usepackage[hidelinks]{hyperref}
\usepackage{natbib}
\usepackage{microtype}
\usepackage{caption}
\usepackage{enumitem}

\bibliographystyle{plainnat}

\title{\textbf{Gradient Boosting Recovers Fill Probability Signal\\
That a Deep Survival Model Misses}\\[0.4em]
\large A false negative about data content, from 70 days of\\
cryptocurrency limit order book data}

\author{Jude Kriel Ramcharitar\\
\small Independent Researcher\\
\small \texttt{j.k.ramcharitar@gmail.com}}

\date{August 2026 \quad$\cdot$\quad Version 3.0 --- Working Paper}

\begin{document}
\maketitle

\begin{abstract}
\noindent
Retail and non-institutional algorithmic traders operating under execution
delays of 200\,ms to 2\,s face substantial stale-quote risk: the book can move
before a quote reaches the exchange. A natural response is to substitute
prediction for speed, estimating the probability that a quote fills given
current limit order book (LOB) state. Whether this is possible at the data
frequency available to such traders is an empirical question, and we find that
whether the available signal is recovered depends strongly on model class.

Using 269{,}330 LOB snapshots collected continuously over 70 days from Kraken
across BTC/USD, ETH/USD and SOL/USD, we benchmark fill probability estimators
against a deliberately trivial reference: a 32-cell lookup table holding only
the mean historical fill rate for each combination of quote offset, latency
window and side. Under a snapshot-grouped temporal split, gradient-boosted trees
over microstructure features beat this reference on all three assets, by
$+0.174$, $+0.097$ and $+0.045$ AUC (mean $+0.105$). An LSTM survival model over
the same features, given a 50-snapshot temporal window the trees never receive,
beats it on none, scoring $-0.012$, $-0.007$ and $-0.116$ against the table.

The positive result survives our leakage and specification checks. It holds under a split
that prevents any snapshot appearing in both training and test data, and it
strengthens when all price-derived columns are removed, ruling out the trees
partitioning on price level as a proxy for time period. A boosted model
restricted to the three conditioning variables reproduces the lookup table to
four decimal places, confirming the comparison is set up correctly.

We also report a conditioning defect that invalidated our own earlier results,
in which quote offset never reached the model, and which we retain as a
reproduction control. Corrected and uncorrected forms differ by 0.15 C-index,
and the uncorrected form performs at chance across five seeds.

The practical implication reverses the framing we began with. Microstructure
signal is present at roughly 90 observed LOB updates per hour and is worth about
0.10 AUC over quote geometry alone. Extracting it, in our experiments, did not
require more data, finer data or a larger model; it required a model class
suited to tabular features. Our tested deep survival architecture, adapted from
the higher-frequency literature, did not recover it.

\vspace{0.8em}
\noindent\textbf{Keywords:} market making, market microstructure, limit order
book, fill probability, execution latency, gradient boosting, model selection

\vspace{0.3em}
\noindent\textbf{JEL:} G12, G14, C45, C61, G23
\end{abstract}

\vspace{1em}

\section{Introduction}

In competitive markets, speed is capital. Firms dominating modern equity and
cryptocurrency liquidity provision have invested heavily in compressing
execution latency to microseconds. Microseconds do not matter to the real
economy; being first to respond to a price signal is profitable, repeatedly, at
scale. This raises a precise question: is speed itself the valuable resource, or
is it a proxy for the ability to act on information before it becomes stale?

If the latter, speed might not be strictly necessary. A participant who cannot
act in microseconds could still act correctly, provided they can predict what
the market will look like when their order arrives. Earlier versions of this
work proposed \emph{latency-compensating market making} on that basis, in which
models trained on LOB microstructure substitute predictive accuracy for
execution speed.

This version reports the empirical status of the predictive component. The
signal exists. Our original model could not find it.

\subsection{Two findings}

\textbf{Microstructure predicts fills at retail data frequency.}
Gradient-boosted trees over microstructure features beat a lookup table on quote
offset, latency and side by $+0.174$, $+0.097$ and $+0.045$ AUC on BTC, SOL and
ETH. The effect appears on every asset, survives a snapshot-grouped split, and
strengthens when price-derived features are removed.

\textbf{A deep survival model extracts none of it.} An LSTM over the same
features, with a 50-snapshot temporal window the trees never see, scores at or
below the lookup table on all three assets. Pooled, it matches the table to
within 0.0003 AUC.

The gap between these findings is the paper. Both models had the same features,
the same labels and the same split. One recovered roughly 0.10 AUC of signal and
the other recovered none.

\subsection{How we arrived here}

We report the sequence because the intermediate conclusions were wrong in
instructive ways.

Version~1 presented results on synthetic LOB data. Those results were
encouraging and uninformative: a geometric Brownian motion price process with a
hand-specified fill rule contains whatever structure the specification puts
there.

On real data our first results fell below chance. A diagnostic located a defect
in the model's input pipeline (Section~\ref{sec:defect}). Correcting it produced
a real improvement, and we then compared the corrected model against a lookup
table, a comparison that should have come first. The model did not beat it. We
concluded microstructure carried no signal at this data frequency and wrote that
conclusion up.

Running a gradient-boosted baseline, which that draft had itself flagged as
missing, overturned the conclusion within twenty minutes. The features were
informative throughout. The measurement instrument was not.

We draw the methodological point out in Section~\ref{sec:lessons}. Inferring a
data limitation from one model's failure is a specific and avoidable error, and
we made it twice, for two different underlying reasons.

\subsection{Contributions}

\begin{enumerate}[itemsep=2pt,topsep=4pt]
\item We show microstructure features predict our replay-defined fill outcome at
roughly 90 observed LOB updates per hour, worth a mean $+0.105$ AUC over quote
geometry, on three assets (Section~\ref{sec:results}).

\item We show an LSTM survival model of the kind adapted from higher-frequency
work extracts none of this signal, despite receiving strictly more information
than the model that does (Section~\ref{sec:comparison}).

\item We release 269{,}330 timestamped 10-level LOB snapshots with the full
collection, feature and evaluation pipeline (Section~\ref{sec:data}).

\item We report a leakage audit for the positive result: grouped splits, price
removal, a shuffled-split contrast and a control fit that reproduces the
reference (Section~\ref{sec:audit}).

\item We document a conditioning defect that invalidated our own prior results
and retain it as a reproduction control (Section~\ref{sec:defect}).
\end{enumerate}

\subsection{Scope}
\label{sec:scope}

Because the interpretive distance between what we measured and what the result
is likely to be read as saying is large, we state the boundaries before the
evidence.

\textbf{Established.} In our Kraken sample, covering three cryptocurrency pairs
from June to August 2026 at roughly 90 observed LOB updates per hour,
microstructure features carry predictive content for our replay-defined fill
outcome beyond quote offset, latency and side, recoverable by gradient-boosted
trees.

\textbf{Not established: economic exploitability.} AUC of 0.78 is not a trading
strategy. We run no backtest, no transaction cost analysis and no live
evaluation. The distance between a predictive edge and a profitable one is where
most of the difficulty in this problem lies, and we do not cross it.

\textbf{Not established: realised fills.} Labels are reconstructed from
subsequent book states and do not model queue position, so the measured outcome
is an optimistic proxy for a late-arriving limit order. We show discrimination
for a replay-defined fill, not for a realised execution.

\textbf{Not established: calibration.} AUC measures ranking. It does not
establish that a predicted 0.70 corresponds to a 70\% fill rate. This matters
because the abstention rule in \eqref{eq:abstain} thresholds a probability, and
a well-ranked but poorly calibrated estimator would not support it without
recalibration.

\textbf{Not established: other deep architectures.} We tested one architecture
at one hyperparameter setting. A fair reading is that our implementation,
hyperparameters or survival objective are at fault rather than deep learning
generally.

\textbf{Not established: generality.} One venue, one 70-day window, three
cryptocurrency pairs, one feed type. We make no claim about equity venues, other
sampling frequencies or other periods.

\textbf{Not tested.} The reinforcement learning agent proposed in Version~1
consumed estimates from the defective model and was trained only on synthetic
data. We report no RL results (Section~\ref{sec:rl-status}).

\section{Related Work}

\subsection{Market making and latency}

The modern framework originates with \citet{garman1976} and \citet{ho1981},
formalised as a stochastic control problem by \citet{avellaneda2008}. Extensions
incorporate adverse selection and execution uncertainty \citep{cartea2018}, and
\citet{cartea2023} introduce stochastic delay into optimal execution.

Nearly all analytical market making models assume negligible execution latency.
\citet{lehalle2022} show the assumption is load-bearing: the value of predictive
signals about future order flow is systematically eroded by latency. Our results
bear on the premise of that argument, since they establish a usable signal
exists at retail data frequency in the first place.

\subsection{The latency arms race}

\citet{moallemi2013} give the first formal model of latency cost.
\citet{budish2015} characterise continuous-time limit order books as generating
structurally wasteful speed investment. \citet{aquilina2022} quantify this,
finding latency arbitrage races roughly once per minute per symbol in FTSE 100
stocks with modal durations of 5--10 microseconds, and estimating that
eliminating them would cut effective spreads by about 17\%.

Proposed remedies (the IEX speed bump of \citealp{hu2019}, frequent batch
auctions in \citealp{budish2015}, micro-burst fees in \citealp{brolley2023})
operate at the exchange level and require regulatory or venue adoption. None
offers recourse to an individual trader on existing infrastructure.

\subsection{Fill probability estimation}

\citet{lo2002} established econometric foundations for execution probability
modelling. \citet{maglaras2022} introduced a deep learning approach treating
execution as a time-to-fill distribution. \citet{arroyo2024} advanced this with
a convolutional-transformer using survival analysis, handling the
right-censoring inherent in limit order data.

That work uses tick-level message data. Our LSTM follows its survival
formulation and fails on snapshot data at a median 18.2-second interval. We do
not read this as contradicting those results. The likelier reading is that
architectures designed for dense sequential data lose their advantage when the
sequence is sparse and irregular, which is a boundary condition worth stating.

\subsection{Tabular deep learning}

Our central comparison connects to a finding from outside finance. Systematic
evaluations report that gradient-boosted trees remain competitive with or
superior to neural networks on tabular data, particularly at moderate sample
sizes with heterogeneous, weakly correlated features
\citep{grinsztajn2022,shwartz2022}. Our features are engineered scalars
(spreads, imbalances, volumes) computed per snapshot, which is tabular in
structure even though the underlying process is sequential. Our result is
consistent with that literature and, so far as we are aware, has not previously
been demonstrated for LOB fill prediction.

\subsection{Reinforcement learning for market making}

\citet{spooner2018} demonstrated RL agents learning market making policies in
simulation; \citet{guo2023} added continuous action spaces; \citet{briola2021}
established viability in cryptocurrency settings. RELAVER \citep{jiang2025}
models random execution delays of 30--100\,ms. We do not extend this line in the
present version (Section~\ref{sec:rl-status}).

\section{Problem Formalisation}
\label{sec:formal}

\subsection{Standard setup}

A market maker operates over horizon $[0,T]$. The mid-price follows
$S_t = S_0 + \sigma W_t$ with $\sigma > 0$ and $W_t$ standard Brownian motion.
Posting a bid at distance $\delta^b$ and ask at $\delta^a$ from mid, fills arrive
as Poisson processes with intensities
\begin{equation}
\lambda^b(\delta^b) = A e^{-k\delta^b}, \qquad
\lambda^a(\delta^a) = A e^{-k\delta^a},
\end{equation}
where $A, k > 0$ govern decay of arrival rate with distance from mid. The
objective is
\begin{equation}
\max_{\delta^b,\delta^a} \; \mathbb{E}\!\left[ X_T + q_T S_T - \gamma q_T^2 \right],
\end{equation}
with $\gamma > 0$ the inventory risk aversion parameter.

\subsection{Execution latency}

We introduce execution latency $\ell$, the delay between quoting decision and
order arrival, and consider $\ell \in [200\,\text{ms}, 2000\,\text{ms}]$. A quote
posted at $t$ arrives at $t + \ell$ into a book state $\mathcal{F}_{t+\ell}$
unobserved at decision time, so the relevant fill probability is
\begin{equation}
P^{\text{fill}}(\delta, \ell, \mathcal{F}_t)
= \mathbb{P}\!\left[\text{order fills} \mid \delta, \mathcal{F}_{t+\ell}\right],
\end{equation}
which depends on a future state. This produces \emph{adverse selection
amplification}, where quotes are picked off as the mid moves during
transmission, and \emph{cancellation impotence}, where an adverse signal cannot
be acted on before a fill occurs.

\subsection{The latency-compensating problem}

Let $\widehat{P}^{\text{fill}}$ estimate fill probability and
$\widehat{S}_{t+\ell}$ the mid-price at arrival. The problem is
\begin{equation}
\max_{\delta^b,\delta^a} \;
\mathbb{E}\!\left[ X_T + q_T \widehat{S}_T - \gamma q_T^2
\;\middle|\; \widehat{\mathcal{F}}_{t+\ell} \right]
\end{equation}
subject to participation constraints
\begin{equation}
\widehat{P}^{\text{fill}}(\delta, \ell, \mathcal{F}_t) \geq \tau_{\text{fill}},
\qquad
\left| \widehat{S}_{t+\ell} - S_t \right| \leq \Delta_{\text{safe}},
\label{eq:abstain}
\end{equation}
with abstention whenever either is violated.

\subsection{Viability conditions}

Version~1 stated three conditions: sufficient predictive accuracy, adequate
uninformed order flow \citep{chen2018}, and strict abstention discipline. This
paper concerns the first. Our result is that a usable estimator
$\widehat{P}^{\text{fill}}$ is obtainable at retail data frequency, which
establishes the condition is satisfiable. Whether the accuracy suffices for
profitability, and whether the other two conditions hold, is untested here.

\section{Data}
\label{sec:data}

\subsection{Collection}

LOB data were collected from Kraken via WebSocket API v2, subscribing to the
\texttt{book} channel at depth 10 for BTC/USD, ETH/USD and SOL/USD. Updates are
timestamped at receipt and stored with the full 10-level ladder in Parquet,
partitioned by asset, date and hour.

Collection ran from 15 June to 24 August 2026. The collector is managed by a
systemd unit with \texttt{Restart=always}; over the final three weeks it
restarted automatically three times after WebSocket disconnections. An earlier
configuration using a detached terminal session and a bounded reconnect loop
failed silently on 20 July and stayed down for fourteen days, costing an
estimated 90{,}000 snapshots. We note this because unattended collection
reliability is a material and under-discussed constraint on independent
microstructure research.

\subsection{Sampling characteristics}

Kraken's book channel is event-driven, so sampling is sparse and irregular.

\begin{table}[ht]
\centering
\caption{Inter-update intervals, BTC/USD.}
\label{tab:gaps}
\begin{tabular}{lr}
\toprule
Statistic & Value \\
\midrule
Mean gap            & 40.1\,s \\
Median gap          & 18.2\,s \\
Minimum gap         & $<$0.01\,s \\
Maximum gap         & 6339\,s \\
Updates per hour    & $\approx$90 \\
\bottomrule
\end{tabular}
\end{table}

\noindent
The maximum gap corresponds to the collector outage rather than market
quiescence. This sampling rate is three to four orders of magnitude coarser than
the tick data used in \citet{arroyo2024} and \citet{maglaras2022}, which is the
context in which our LSTM result should be read.

\begin{table}[ht]
\centering
\caption{Collected dataset, 15 June -- 24 August 2026.}
\label{tab:data}
\begin{tabular}{lrrr}
\toprule
Asset & Snapshots & Label rows & Overall fill rate \\
\midrule
BTC/USD & 111{,}002 & 3{,}551{,}392 & 0.734 \\
ETH/USD &  34{,}380 & 1{,}099{,}488 & 0.754 \\
SOL/USD & 123{,}948 & 3{,}965{,}664 & 0.534 \\
\midrule
Total   & 269{,}330 & 8{,}616{,}544 & --- \\
\bottomrule
\end{tabular}
\end{table}

\subsection{Label construction}

For each snapshot we simulate limit order submissions at offsets
$\delta \in \{1,2,3,5\}$ ticks on each side, recording whether each would have
filled within each latency window $\ell \in \{200,500,1000,2000\}$\,ms based on
subsequent book evolution. This yields $4 \times 4 \times 2 = 32$ label rows per
snapshot, with unfilled orders right-censored at $\ell$.

Both conditioning variables show strong monotone structure
(Table~\ref{tab:marginals}), which is what makes the lookup table a demanding
reference rather than a straw one.

\begin{table}[ht]
\centering
\caption{Marginal fill rates by quote offset and latency window, SOL/USD.}
\label{tab:marginals}
\begin{tabular}{lr@{\hskip 2.5em}lr}
\toprule
Offset (ticks) & Fill rate & Latency (ms) & Fill rate \\
\midrule
1 & 0.782 & 200  & 0.385 \\
2 & 0.550 & 500  & 0.498 \\
3 & 0.446 & 1000 & 0.585 \\
5 & 0.358 & 2000 & 0.669 \\
\bottomrule
\end{tabular}
\end{table}

\subsection{Features}

From each snapshot we construct 109 features: 40 raw price and volume levels,
normalised prices and volumes, mid-price, spread, spread in basis points, order
book imbalance, depth imbalance, and rolling statistics and lagged differences.
Prices are normalised by contemporaneous mid, volumes by a rolling mean.

We define a \emph{microstructure} subset of 63 features by removing every
price-derived column, raw and normalised, retaining spreads, imbalances, order
book imbalance, depth, volumes and volatility measures. This subset is used for
the headline results; the rationale is in Section~\ref{sec:price}.

Two pipeline parameters required adjustment for sparse data. The rolling
normalisation window was reduced from 1000 to 10 snapshots and lag features from
$\{1,5,10,20,50\}$ to $\{1,2,3,5\}$. At the original settings a 1000-snapshot
window spans roughly eleven hours of wall-clock time, which is not a meaningful
local normalisation.

\section{Models and Evaluation}
\label{sec:models}

\subsection{The reference}

Our benchmark is a lookup table: the mean fill rate for each
$(\delta, \ell, \text{side})$ cell, 32 cells, estimated on training data and
applied unchanged to test data. It has no learned parameters, no access to the
order book and no temporal information.

It is a demanding reference. Quote offset spans a 0.42 range in fill rate and
latency 0.28 (Table~\ref{tab:marginals}). Any model failing to beat it adds
nothing over two variables the trader selects directly and already knows.

\subsection{Gradient-boosted trees}

We fit gradient-boosted trees (XGBoost, 400 estimators, depth 6, learning rate
0.05, subsample 0.8) on the microstructure subset for the current snapshot only,
plus the three conditioning variables. The model receives \textbf{no temporal
window}.

\subsection{LSTM survival model}

The LSTM encodes a lookback of $L = 50$ snapshots with two layers (hidden
dimension 64, dropout 0.2). A conditioning vector
$c = [\ell/\ell_{\max},\, \delta/\delta_{\max},\, \text{side}] \in \mathbb{R}^3$
passes through an MLP and is concatenated before a monotone survival decoder.
Total parameters: 90{,}754. Training minimises the negative log-likelihood under
right-censoring,
\begin{equation}
\mathcal{L} = -\frac{1}{N}\sum_{i=1}^{N}
\left[ \Delta_i \log h(\widetilde{T}_i) + \log S(\widetilde{T}_i) \right],
\end{equation}
using Adam at $\eta = 10^{-3}$ with early stopping.

The comparison is deliberately unfavourable to the trees. The LSTM sees fifty
snapshots of history and 46 additional price columns; the trees see one snapshot
and fewer features.

\subsection{Splitting}
\label{sec:split}

All results use a \emph{snapshot-grouped} temporal split at 70/15/15. Because
each snapshot generates 32 label rows sharing identical features, a row-wise
split lets a snapshot appear in both training and test data, allowing a
sufficiently flexible model to memorise it. Grouping on snapshot removes this.
The effect is measured in Section~\ref{sec:audit}.

\subsection{Metric}

We report fill AUC: area under the ROC curve for the binary filled/not-filled
task, computed exactly via the Mann-Whitney rank-sum identity with ties credited
0.5. AUC is the relevant quantity because the abstention rule in
\eqref{eq:abstain} tests $\widehat{P}^{\text{fill}} \geq \tau_{\text{fill}}$, a
binary decision that never compares timings.

\section{A Conditioning Defect and Its Use as a Control}
\label{sec:defect}

Our first real-data results were substantially worse than chance. Our initial
hypothesis was insufficient data. That hypothesis was wrong, and the diagnostic
that displaced it is worth describing.

\subsection{The defect}

The conditioning vector $c$ originally contained only $\ell$. Quote offset
$\delta$ and side were used to construct labels but never reached the model.

For a fixed snapshot and latency, the eight label rows spanning four offsets and
two sides therefore present \emph{identical} model input while carrying
different targets, with marginal fill rates from 0.782 at 1 tick to 0.358 at 5
ticks. No function of the input can separate them, and the loss-minimising
response is to predict the conditional mean and ignore the input. Validation
loss moved roughly 0.003 over 40 epochs in every run.

This also explains an observation that misled us. Across collection milestones,
measured C-index moved 0.5008, 0.5308, 0.5526 as the dataset grew from 21K to
51K to 103K snapshots. We read this as convergence. All three lie within seed
variance of chance, and a later run on 269K snapshots returned 0.4749, worse
with more data. That is the expected behaviour when additional data supplies
additional mutually contradictory input-target pairs.

\subsection{Correction and control}

Adding $\delta$ and side to $c$ raised C-index from 0.4749 to 0.6266 on an
otherwise identical configuration, with validation loss movement increasing
roughly eightfold.

We retain the defective configuration as an arm, \textsc{lob}, in which the LOB
window is supplied but $\delta$ and side are reset to constants. It is a
reproduction control: any claim that a result derives from LOB features should
come with a demonstration that the corresponding defective configuration
performs at chance. Across five seeds it scores $0.5034 \pm 0.031$
(Table~\ref{tab:ablation}).

\subsection{Consequence for prior results}

All real-data figures in Versions 1.0--2.0 were produced by the defective
configuration and should be disregarded, including the C-index progression above
and a projection that 300{,}000 snapshots would be needed for convergence, which
was extrapolated from noise.

\section{Results}
\label{sec:results}

\subsection{Main comparison}
\label{sec:comparison}

Table~\ref{tab:main} reports all three predictors on identical test rows under
the grouped split.

\begin{table}[ht]
\centering
\caption{Fill AUC by model class, grouped temporal split, identical test rows.
GBM uses the 63-feature microstructure subset and the current snapshot only; the
LSTM uses all 109 features and a 50-snapshot window.}
\label{tab:main}
\small
\begin{tabular}{lrrrrrr}
\toprule
Asset & Test rows & Lookup & GBM & LSTM & GBM $-$ Lookup & LSTM $-$ Lookup \\
\midrule
BTC/USD & 29{,}788 & 0.6390 & 0.8128 & 0.6271 & $+0.1738$ & $-0.0119$ \\
SOL/USD & 29{,}845 & 0.6819 & 0.7786 & 0.6750 & $+0.0967$ & $-0.0069$ \\
ETH/USD & 29{,}798 & 0.6505 & 0.6956 & 0.5341 & $+0.0450$ & $-0.1164$ \\
\midrule
Mean    &          &        &        &        & $+0.1052$ & $-0.0451$ \\
\bottomrule
\end{tabular}
\end{table}

\noindent
Three observations.

The GBM beats the lookup table on every asset, by $+0.174$, $+0.097$ and
$+0.045$. Microstructure features carry predictive content for the replay-defined
fill outcome at this sampling rate.

The LSTM beats it on none, and on ETH/USD underperforms by 0.116. Pooled across
assets it matched the table to $+0.0003$ AUC.

The effect size varies fourfold across assets and does not track book thickness
in any obvious way. SOL/USD has the thinnest book and the steepest offset
gradient but a middling gap; ETH/USD has both the fewest snapshots and the
smallest gap, which may reflect its 34{,}380 snapshots rather than any property
of the asset. We report the range rather than leading with the mean.

\subsection{The LSTM failure}

The LSTM receives strictly more information than the GBM: the same
microstructure features, plus a 50-snapshot window, plus 46 additional
price-derived columns. It performs worse on all three assets.

Its training behaviour is consistent with underfitting rather than overfitting.
On ETH it stopped at epoch 13 with validation loss moving 1.7405 to 1.7329; on
SOL at epoch 12; on BTC at epoch 32 with the largest movement, 1.6592 to 1.6183,
and the best relative result. Longer patience, a larger hidden dimension or a
different learning rate schedule might change this, and we did not run that
search.

What we can claim is narrower than ``deep learning fails here.'' A specific
architecture, adapted from work on tick-level data and trained with our
hyperparameters, does not recover signal that a standard tabular method recovers
easily from less information.

\subsection{Ablation}
\label{sec:ablation}

For completeness we report the ablation run on the LSTM before the GBM
comparison existed, using C-index on a pooled row-wise split.

\begin{table}[ht]
\centering
\caption{LSTM ablation, C-index, five seeds per arm, pooled row-wise split.}
\label{tab:ablation}
\begin{tabular}{lrrrr}
\toprule
Arm & Mean & SD & Min & Max \\
\midrule
\textsc{full} (LOB $+$ conditioning) & 0.6101 & 0.034 & 0.5670 & 0.6536 \\
\textsc{cond} (conditioning only)    & 0.5247 & 0.041 & 0.4853 & 0.5885 \\
\textsc{lob}  (defective control)    & 0.5034 & 0.031 & 0.4678 & 0.5427 \\
\bottomrule
\end{tabular}
\end{table}

\noindent
\textsc{full} exceeds \textsc{cond} by 0.085, which taken alone appears to show
the LSTM using LOB features productively. It does not. \textsc{cond} is a model
trained with its dominant input stream zeroed, a damaged network rather than a
fair alternative predictor, and Table~\ref{tab:main} shows the LSTM adds nothing
over quote geometry. We report this because the ablation was run first and,
alone, supported a conclusion the external reference reversed.

\textsc{lob} at 0.5034 confirms the defect of Section~\ref{sec:defect}
reproduces to chance.

\subsection{Status of the reinforcement learning component}
\label{sec:rl-status}

Version~1 proposed a latency-aware PPO agent consuming fill probability
estimates, with a reward penalising adverse fills attributable to execution
delay. We report no RL results. The agent was trained only on synthetic data,
where it did not converge, and it consumed estimates from the defective model,
which makes those runs uninterpretable.

Given the results above, an RL agent is now worth revisiting with GBM-derived
fill estimates. We regard this as the most promising immediate extension and it
is not attempted here.

\section{Audit of the Positive Result}
\label{sec:audit}

A positive result obtained after several negative ones deserves scrutiny. We
report the checks we ran, including those that failed to find a problem.

\subsection{Splitting}

Our initial GBM result used a row-wise temporal split and scored 0.8264 on
BTC/USD. Because 32 label rows share each snapshot's features, this permits a
snapshot to straddle train and test. Regrouping on snapshot reduced the score to
0.8127, a loss of 0.014. The effect is real but small, and all results in
Section~\ref{sec:results} use the grouped split.

\subsection{Shuffled-split contrast}

A shuffled split, invalid for time series, scored 0.8258 against 0.8264 for the
temporal row split. The near-identity suggests the signal is local and
stationary rather than trend-dependent. We note this as an observation rather
than a check that was passed, since we do not have a clean interpretation of it.

\subsection{Price-level features}
\label{sec:price}

Feature importances on BTC initially concerned us. The largest contributors were
\texttt{bid\_price\_2} (0.143) and \texttt{ask\_price\_2} (0.119), raw prices in
the tens of thousands of dollars. A tree splitting on absolute price level in
temporally-ordered data is partitioning on time period, which proxies for
regime, and that is not microstructure prediction.

Removing every price-derived column, raw and normalised, tests this directly.

\begin{table}[ht]
\centering
\caption{Fill AUC gain over the lookup table by feature set, grouped split.}
\label{tab:price}
\begin{tabular}{lrrrr}
\toprule
Feature set & LOB features & BTC & SOL & ETH \\
\midrule
All                 & 109 & $+0.1658$ & $+0.0926$ & $+0.0233$ \\
No price columns    &  68 & $+0.1724$ & $+0.0982$ & $+0.0461$ \\
Microstructure only &  63 & $+0.1738$ & $+0.0967$ & $+0.0450$ \\
Conditioning only   &   0 & $-0.0007$ & $+0.0000$ & $-0.0007$ \\
\bottomrule
\end{tabular}
\end{table}

\noindent
Removing price columns \emph{improved} every asset, by $+0.008$ on BTC,
$+0.022$ on ETH and $+0.006$ on SOL. Our hypothesis was wrong: the price columns
were nuisance variables consuming splits, not a leakage channel. Headline
figures use the 63-feature microstructure subset accordingly.

\subsection{Validating the comparison}

The conditioning-only row of Table~\ref{tab:price} fits a GBM on the three
conditioning variables alone. It reproduces the lookup table to within 0.0007 on
every asset. That is the expected behaviour, and it confirms the comparison
measures what we intend.

\subsection{What remains unchecked}

We have not tested other venues, other feed types or other tree
hyperparameters. We have not run repeated seeds for the GBM, so per-asset
figures carry unmeasured variance, though the between-asset spread ($+0.045$ to
$+0.174$) is larger than seed variance plausibly accounts for. We have not
established that the signal is economically exploitable.

\section{Methodological Lessons}
\label{sec:lessons}

Three errors in this project were caught only by checks that could have been run
at the outset, at a combined cost of about two months. Each is cheap to avoid.

\textbf{Run the trivial reference first.} The lookup table took minutes to
implement and reversed our conclusion twice: first by showing the corrected LSTM
added nothing, then by providing the scale against which the GBM's gain was
measurable. Without it we would have published the LSTM's 0.6266 C-index as a
modest positive result.

\textbf{An ablation is not a substitute for an external baseline.} Our LSTM ablation
showed a clean $+0.085$ separation between \textsc{full} and \textsc{cond} and
would, presented alone, have supported the claim that the model used LOB
features productively. Both arms were below the lookup table. Ablations measure
internal dependence, not external validity.

\textbf{Do not infer a data limitation from one model's failure.} We reached
that inference twice, for different reasons, and it is worth separating them.

The first time, a conditioning defect was feeding the model contradictory
input-target pairs (Section~\ref{sec:defect}). Its flat C-index looked like a
sample-size problem, and we projected that 300{,}000 snapshots would be needed
for convergence. We then collected for six weeks, adding roughly 130{,}000
snapshots, and the projection was extrapolated from noise the whole time. More
data cannot repair a defective input pipeline.

The second time, with the defect corrected, the model matched a lookup table and
we concluded the data contained no usable signal. A different model class
overturned that in twenty minutes.

Neither episode licenses the opposite rule of blaming the model first. Failures
are sometimes genuinely about data, labels, optimisation or implementation. The
transferable version is narrower: before concluding that data contain no usable
signal, rule out implementation defects and test at least one model with a
different inductive bias.

\section{Limitations}

\textbf{No economic evaluation.} AUC is not profit. We run no backtest, model no
transaction costs, and make no claim that the signal survives the bid-ask
spread, maker fees, queue position or adverse selection. This is the largest gap
between our result and its motivating application.

\textbf{Simulated fills.} Labels come from replaying subsequent book states, not
executed orders. Queue position is not modelled, so we assume a quote at a price
level fills when that level trades through, which is optimistic for a latecomer.

\textbf{One venue.} Kraken's depth-10 book channel only.

\textbf{One deep architecture.} We tested an LSTM survival model with one
hyperparameter configuration. The convolutional-transformer of
\citet{arroyo2024}, which Version~1 proposed, was not re-evaluated after the
conditioning fix.

\textbf{No GBM seed variance.} Per-asset GBM figures come from single fits.

\textbf{ETH sample size.} ETH/USD has 34{,}380 snapshots against 111{,}002 and
123{,}948 for the others, and shows both the weakest GBM gain and the worst LSTM
failure. Its results should be read as confounded with sample size.

\section{Conclusion}

We asked whether limit order book microstructure predicts fill probability at
the data frequency available to a retail participant, roughly 90 updates per
hour from a public exchange feed. For our replay-defined fill outcome, it does.
Gradient-boosted trees over
microstructure features beat a lookup table on quote offset, latency and side by
$+0.174$, $+0.097$ and $+0.045$ AUC across three cryptocurrency pairs, under a
snapshot-grouped temporal split, with price-derived features removed, and with a
conditioning-only control that reproduces the reference exactly.

An LSTM survival model over the same features, given a fifty-snapshot temporal
window the trees never receive, recovers none of it. It matches the trivial
reference pooled and underperforms it on all three assets individually.

We set out to test whether prediction can substitute for speed. We can now say
the predictive component is available at retail data frequency, which was
genuinely in doubt, and that in our experiments recovering it turned on model
selection rather than on collecting more data. Whether the resulting edge
survives transaction costs, queue position and adverse selection is the next
question and the harder one.

Our own route to this result ran through the opposite conclusion, written up in
full. The correction came from a twenty-minute reference we had deferred for two
months. We report that sequence because the cost of the error was six weeks of
data collection aimed at a problem that did not exist.

\section*{Data and Code Availability}
\addcontentsline{toc}{section}{Data and Code Availability}

The complete collection pipeline, feature engineering, models and every
diagnostic reported here are available at
\url{https://github.com/azure-cyber/latency-compensating-mm}. The scripts
underlying Sections~\ref{sec:results} and \ref{sec:audit}
(\texttt{gbm\_matched.py}, \texttt{gbm\_price\_test.py},
\texttt{gbm\_leakage\_check.py}, \texttt{ablation.py}, \texttt{diagnose.py}) are
included so both the positive and negative results can be checked directly.

\section*{Disclosure of AI Assistance}
\addcontentsline{toc}{section}{Disclosure of AI Assistance}

The author used Claude (Anthropic) extensively throughout this project: for
implementation of the data collection pipeline, feature engineering, models and
evaluation code; for diagnostic analysis including identification of the
conditioning defect in Section~\ref{sec:defect} and design of the audit in
Section~\ref{sec:audit}; and for drafting and revision of this manuscript. The
author designed the study, directed the analysis, made all substantive research
decisions, verified the results, and is solely responsible for the content and
any errors.

\section*{Acknowledgements}
\addcontentsline{toc}{section}{Acknowledgements}

This work was conducted independently, without institutional affiliation,
funding or supervision. The author declares no conflicts of interest and holds
no positions in any asset discussed.

\bibliography{refs}

\end{document}
